% from ega0-1-fr.tex
\oldpage[0\textsubscript{I}]{62}

\begin{corollary}[7.1.7]
\label{0.7.1.7-fr}
Un anneau préadmissible noethérien admet un plus grand idéal de définition.
\end{corollary}

\begin{corollary}[7.1.8]
\label{0.7.1.8-fr}
Si un anneau préadmissible $A$ est tel que, pour un idéal de définition
$\mathfrak J$, les puissances $\mathfrak J^n$ ($n>0$) forment un système
fondamental de voisinages de $0$, il en est de même des puissances
${\mathfrak J'}^n$ de tout idéal de définition $\mathfrak J'$ de $A$.
\end{corollary}

\begin{definition}[7.1.9]
\label{0.7.1.9-fr}
On dit qu'un anneau préadmissible $A$ est préadique s'il existe un idéal de
définition $\mathfrak J$ de $A$ tel que les $\mathfrak J^n$ forment un système
fondamental de voisinages de $0$ dans $A$ (ou, ce qui revient au même, tel que
les $\mathfrak J^n$ soient ouverts). On appelle anneau adique un anneau
préadique séparé et complet.
\end{definition}

Si $\mathfrak J$ est un idéal de définition d'un anneau préadique (resp.
adique) $A$, on dit encore que $A$ est un anneau
\emph{$\mathfrak J$-préadique} (resp. \emph{$\mathfrak J$-adique}), et que sa
topologie est la topologie \emph{$\mathfrak J$-préadique} (resp.
\emph{$\mathfrak J$-adique}). Plus généralement, si $M$ est un $A$-module, la
topologie sur $M$ ayant pour système fondamental de voisinages de $0$ les
sous-modules $\mathfrak J^nM$ est dite topologie
\emph{$\mathfrak J$-préadique} (resp. \emph{$\mathfrak J$-adique}). En vertu
de (7.1.8), ces topologies sont indépendantes du choix de l'idéal de définition
$\mathfrak J$.

\begin{proposition}[7.1.10]
\label{0.7.1.10-fr}
Soient $A$ un anneau admissible, $\mathfrak J$ un idéal de définition de $A$.
Alors $\mathfrak J$ est contenu dans le radical de $A$.
\end{proposition}

Cet énoncé est équivalent à l'un quelconque des corollaires suivants~:

\begin{corollary}[7.1.11]
\label{0.7.1.11-fr}
Pour tout $x\in\mathfrak J$, $1+x$ est inversible dans $A$.
\end{corollary}

\begin{corollary}[7.1.12]
\label{0.7.1.12-fr}
Pour que $f\in A$ soit inversible dans $A$, il faut et il suffit que son image
canonique dans $A/\mathfrak J$ soit inversible dans $A/\mathfrak J$.
\end{corollary}

\begin{corollary}[7.1.13]
\label{0.7.1.13-fr}
Pour tout $A$-module $M$ de type fini, la relation $M=\mathfrak JM$
(équivalente à $M\otimes_A(A/\mathfrak J)=0$) entraîne $M=0$.
\end{corollary}

\begin{corollary}[7.1.14]
\label{0.7.1.14-fr}
Soit $u:M\to N$ un homomorphisme de $A$-modules, $N$ étant de type fini~;
pour que $u$ soit surjectif, il faut et il suffit que
$u\otimes1:M\otimes_A(A/\mathfrak J)\to N\otimes_A(A/\mathfrak J)$ le soit.
\end{corollary}

En effet, l'équivalence de (7.1.10) et (7.1.11) résulte de Bourbaki,
\emph{Alg.}, chap.~VIII, \textsection6, n\textsuperscript{o}~3, th.~1, et
celle de (7.1.10) et (7.1.13) de \emph{loc. cit.}, th.~2~; le fait que
(7.1.10) entraîne (7.1.14) résulte de \emph{loc. cit.}, cor.~4 de la prop.~6~;
d'autre part, (7.1.14) entraîne (7.1.13) en l'appliquant à l'homomorphisme nul.
Enfin, (7.1.10) entraîne que si $f$ est inversible dans $A/\mathfrak J$, $f$
n'est contenu dans aucun idéal maximal de $A$, donc $f$ est inversible dans
$A$, autrement dit (7.1.10) entraîne (7.1.12)~; et inversement, (7.1.12)
entraîne (7.1.11).

Tout revient donc à démontrer (7.1.11). Or, comme $A$ est séparé et complet et
que la suite $(\mathfrak J^n)$ tend vers $0$, il est immédiat que la série
$\sum_{n=0}^{\infty}(-1)^n x^n$ est convergente dans $A$ et que, si $y$ est
sa somme, on a $y(1+x)=1$.

\subsection{Anneaux adiques et limites projectives}
\label{subsection:0.7.2-fr}

\begin{env}[7.2.1]
\label{0.7.2.1-fr}
Toute limite projective d'anneaux discrets est évidemment un anneau
linéairement topologisé, séparé et complet. Inversement, soit $A$ un anneau
linéairement topologisé, et soit $(\mathfrak J_\lambda)$ un système fondamental
de voisinages ouverts de $0$ dans $A$ formé

